\Needspace{8\baselineskip}\originalchapter{环形力平衡与高 beta 模型}\label{ch:highbeta}
\originalsection{为什么环形柱会向外移动}
把直柱弯成环, 外侧环向磁场较弱, 磁压和箍力不再对称; 内压及极向磁场也产生净向外的环向膨胀力. 外加竖直磁场与环向电流的洛伦兹力可抵消它. 在大长径比圆截面模型, 常写所需场大小
\begin{equation}
 |B_v|\simeq\frac{\mu_0I}{4\pi R_0}
 \left[\ln\frac{8R_0}a+\beta_p+\frac{l_i}2-\frac32\right],
 \quad \beta_p=\frac{2\mu_0\langle p\rangle}{B_\theta(a)^2}.
\end{equation}
$l_i$ 依电流剖面定义内部电感, 符号由电流方向定. 该式是细环渐近结果, 不是任意椭圆或紧凑装置的精确公式. 移动磁轴的 Shafranov 位移必须在整体场方程中计算, 不能靠把压力项加到单粒子漂移得到.

\originalsection{小参数的重新排序}
设 $\epsilon=a/R_0\ll1$, 常规低压展开把 beta 当更小量. 当 $\beta_t=2\mu_0\langle p\rangle/B_0^2\sim\epsilon$, 位移可能与小半径同阶, 原来的微小位移展开不一致. 高 beta tokamak (HBT) 模型改用磁面标签和位移, 保留主要压力与曲率耦合. 它提高某些环形效应的阶数, 仍不包含所有有限长径比修正.

课程的简化剖面以 $\rho=r/a$ 和参数 $\nu=\beta_tq_*^2/\epsilon$ 表示. 边缘绕行积分可得 $q_a=q_*/\sqrt{1-\nu^2}$. 固定 I (即固定 q*) 加热时 \ensuremath{\nu} 增大, \ensuremath{\nu}\ensuremath{\to}1导致 qa 发散并接近分离面. 这个结果同时说明理想平衡存在限制和 q 的定义为何需要沿整条面积分. 大压力并非总能以同一固定电流平衡吸收.

\originalsection{位移、形变和壁之间的区别}
磁面偏移可改变最小 R 与局部场强, 拉长 \ensuremath{\kappa} 可改变电流与 q 的关系, 三角度 \ensuremath{\delta} 则改变曲率和磁面形状. 三种几何效应不能由同一个``等效半径''替代. 即便得到嵌套平衡磁面, 仍须检查理想能量、外部壁和电阻时间尺度. 本册期中题用有限参数矩投影实际求形变, 将近似方程残差与数值求积误差分别报告.
